\section{Final Test Evaluation and Error Analysis}\label{sec:final}
\subsection{Chronological test protocol and selection rule}
Notebook 07 refits the h24 estimators on all \num{7825} eligible development hours through 30 November and evaluates 720 hours from 1--30 December. The h1 development set contains \num{7848} hours and uses the same December timestamps. Model weights remain fixed during each evaluation, while historical-load predictors update with measurements from completed hours or days. These are updating forecasts, not an open-loop prediction of December from a single November cutoff. Observed target-hour weather remains an oracle input.

The final identity follows minimum development mean CV RMSE: absolute-target SVR for both settings. For h24, the notebook evaluates all default, tuned absolute and tuned residual variants; for h1, it evaluates each family's CV-preferred tuned formulation. The h1 table therefore contains fewer ML variants and does not establish a December ranking of all ten searched h1 formulations.

December is held outside the training and model-selection folds, but the repository does not establish a strictly single-use test. Whole-year descriptive EDA includes it, and the progress documentation records repeated final evaluation for reproducibility checks. The notebook's selection expression reads CV scores, even though it appears after calculation of the test table. This supports a CV-based selection rule, rather than prospective proof that every human modeling decision preceded all exposure to December.

\subsection{h24 performance and the residual-SVR result}
\Cref{tab:test24} reports the complete saved h24 comparison. The selected SVR achieves RMSE \num{1624.6} kW, MAE \num{1094.6} kW, MAPE \num{1.78}\% and R$^2=\num{0.9869}$. Against Naive-24h at \num{3296.3} kW RMSE, the reduction is \num{50.71}\%. Naive-168h is stronger than Naive-24h in December, with RMSE \num{2619.8} kW, but remains worse than the selected SVR. Thus the model improves on both seasonal reference rules in this period.
\begin{table}[H]\centering\small
\caption{December test, h24 (T08), ordered by test RMSE. The model chosen on CV before the test is SVR-tuned. RMSE/MAE in kW; skill is the RMSE reduction against Naive-24h.}\label{tab:test24}
\setlength{\tabcolsep}{4pt}
\begin{tabular}{lrrrrr}
\toprule
Model & RMSE & MAE & MAPE (\%) & R$^2$ & Skill (\%) \\
\midrule
SVR $\Delta$-tuned & \num{1419.1} & \num{1046.1} & \num{1.72} & \num{0.9900} & \num{56.95} \\
SVR-tuned & \num{1624.6} & \num{1094.6} & \num{1.78} & \num{0.9869} & \num{50.71} \\
XGB-tuned & \num{1711.5} & \num{1233.5} & \num{2.03} & \num{0.9854} & \num{48.08} \\
XGB $\Delta$-tuned & \num{1719.2} & \num{1209.9} & \num{2.00} & \num{0.9853} & \num{47.84} \\
RF $\Delta$-tuned & \num{1761.9} & \num{1171.4} & \num{1.93} & \num{0.9846} & \num{46.55} \\
XGB-default & \num{1891.0} & \num{1330.3} & \num{2.15} & \num{0.9822} & \num{42.63} \\
DT $\Delta$-tuned & \num{2005.2} & \num{1316.7} & \num{2.20} & \num{0.9800} & \num{39.17} \\
RF-tuned & \num{2017.1} & \num{1331.7} & \num{2.19} & \num{0.9798} & \num{38.81} \\
RF-default & \num{2044.3} & \num{1360.8} & \num{2.24} & \num{0.9792} & \num{37.98} \\
SVR-default & \num{2063.7} & \num{1300.5} & \num{2.20} & \num{0.9788} & \num{37.39} \\
LR-tuned & \num{2169.6} & \num{1624.8} & \num{2.74} & \num{0.9766} & \num{34.18} \\
LR-default & \num{2169.6} & \num{1624.8} & \num{2.74} & \num{0.9766} & \num{34.18} \\
LR $\Delta$-tuned & \num{2217.3} & \num{1629.1} & \num{2.74} & \num{0.9755} & \num{32.73} \\
Naive-168h & \num{2619.8} & \num{1775.4} & \num{2.87} & \num{0.9658} & \num{20.52} \\
DT-tuned & \num{2846.6} & \num{1960.3} & \num{3.19} & \num{0.9597} & \num{13.64} \\
DT-default & \num{2989.6} & \num{1987.4} & \num{3.28} & \num{0.9555} & \num{9.30} \\
Naive-24h & \num{3296.3} & \num{2071.5} & \num{3.39} & \num{0.9459} & \num{0.00} \\
\bottomrule\end{tabular}
\end{table}


Residual SVR has the lowest observed test RMSE, \num{1419.1} kW, with MAE \num{1046.1} kW, MAPE \num{1.72}\% and R$^2=\num{0.9900}$. It was slightly worse on development CV: \num{2427.1} versus \num{2410.6} kW. Close development scores can reverse on a different temporal regime, particularly when error is concentrated on a small number of dates. The saved error analysis identifies the early-December holiday period as a weakness of absolute SVR. Carrying forward the daily load level and learning a correction is a plausible reason for the residual model's different behavior; it is not a demonstrated causal explanation of the entire test advantage. Parameters also differ between the two SVR formulations.

The residual result is reported transparently without promoting that variant to the final selection. Doing so after inspecting December would use test evidence for a new selection decision. Similarly, absolute XGB has slightly lower test RMSE than residual XGB, despite the residual version's better development score. Residual targets therefore do not improve every family in every period.

The selected SVR's December RMSE is approximately \num{32.60}\% below its mean development CV RMSE. December contains a different load regime and the final model has more training observations than earlier folds. This difference is consistent with an easier test period and a larger training set; it does not establish that CV estimates are universally conservative or that all models improve by the same proportion.

\subsection{h1 test results}
\Cref{tab:test1} gives the separately tuned hourly evaluation. The selected absolute SVR reaches RMSE \num{1195.8} kW, MAE \num{861.8} kW, MAPE \num{1.41}\% and R$^2=\num{0.9929}$. Its RMSE reductions are \num{63.72}\% against Naive-24h and \num{78.16}\% against hourly persistence. These percentages use the same 720-hour December target population within h1.
\begin{table}[tbp]\centering\small
\caption{December h1 evaluation (T08), ordered by observed RMSE. The CV-selected model is SVR-tuned. This ordering does not redefine selection. RMSE/MAE in kW; skill is RMSE reduction against Naive-24h.}\label{tab:test1}
\setlength{\tabcolsep}{4pt}
\begin{tabular}{lrrrrr}
\toprule
Model & RMSE & MAE & MAPE (\%) & R$^2$ & Skill (\%) \\
\midrule
SVR-tuned & \num{1195.8} & \num{861.8} & \num{1.41} & \num{0.9929} & \num{63.72} \\
XGB $\Delta$-tuned & \num{1256.6} & \num{826.6} & \num{1.37} & \num{0.9921} & \num{61.88} \\
RF $\Delta$-tuned & \num{1394.1} & \num{918.4} & \num{1.53} & \num{0.9903} & \num{57.71} \\
LR $\Delta$-tuned & \num{1997.2} & \num{1481.5} & \num{2.42} & \num{0.9801} & \num{39.41} \\
DT $\Delta$-tuned & \num{2595.1} & \num{1458.0} & \num{2.46} & \num{0.9665} & \num{21.27} \\
Naive-168h & \num{2619.8} & \num{1775.4} & \num{2.87} & \num{0.9658} & \num{20.52} \\
Naive-24h & \num{3296.3} & \num{2071.5} & \num{3.39} & \num{0.9459} & \num{0.00} \\
Persistence & \num{5475.3} & \num{3946.6} & \num{6.23} & \num{0.8508} & \num{-66.11} \\
MeanProfile & \num{9581.7} & \num{8947.1} & \num{14.63} & \num{0.5431} & \num{-190.68} \\
\bottomrule\end{tabular}
\end{table}


SVR has the lowest RMSE among the h1 models evaluated on test, whereas residual XGB has lower MAE and MAPE: \num{826.6} kW and \num{1.37}\%. There is consequently no metric-independent winner. The h1 individual forecast series was not saved, so the reported metrics cannot be independently reconstructed from stored predictions without refitting. They remain traceable to the test CSV and notebook output. By contrast, all 18 saved h24 prediction columns were checked numerically against their reported metrics during the research audit.

\subsection{Temporal errors and failure cases}
\Cref{fig:test_series} shows the selected h24 SVR following the dominant daily cycle, with a zoom on 11--17 December. Close trajectories over much of the month coexist with visible local errors. The residual plot in \cref{fig:test_residuals} concentrates the largest negative deviations near the beginning of December. With error defined as prediction minus actual, mean bias is $-\num{202.1}$ kW, while individual errors range from $-\num{9362.1}$ to $+\num{6514.2}$ kW. A near-zero median or high R$^2$ does not imply symmetric tails, independent errors or absence of operationally relevant failures.
\begin{figure}[tbp]\centering
\includegraphics[width=\linewidth]{F17_test_actual_vs_pred.png}
\caption{Actual and selected h24 SVR test forecasts, with an 11--17 December zoom and the daily naive reference (F17). Measurements from completed prior days update the load-history inputs throughout the month.}\label{fig:test_series}
\end{figure}
\begin{figure}[tbp]\centering
\includegraphics[width=\linewidth]{F17b_residuals.png}
\caption{Selected h24 SVR errors over December and their empirical distribution (F17b). Negative errors indicate underprediction. The early-month negative tail warrants attention despite the small average bias.}\label{fig:test_residuals}
\end{figure}

The mean December load is \num{63750.8} kW. MAE divided by this mean is about \num{1.72}\%, which differs from MAPE, \num{1.78}\%, because MAPE divides each error by its own observation before averaging. \Cref{fig:test_groups} reveals that improvements over Naive-24h vary across groups. Nonholiday Monday MAE is approximately \num{1007} versus \num{5294} kW, whereas Wednesday MAE is approximately \num{710} versus \num{562} kW: the selected model is not better in every subgroup.
\begin{figure}[tbp]\centering
\includegraphics[width=\linewidth]{F18_error_by_hour_weekday.png}
\caption{December MAE by hour and weekday for selected h24 SVR and Naive-24h (F18). Holidays are separated from weekdays; the holiday group contains only 1 December, limiting generalization of that comparison.}\label{fig:test_groups}
\end{figure}

\Cref{tab:worst_days} identifies 1 December, marked as Mawlid in the implemented calendar, as the largest daily-error case: MAE \num{4579.2} kW and bias $-\num{3616.7}$ kW. The forecast daily mean is \num{54957.4} kW against actual \num{58574.2} kW. A shared religious-holiday indicator learned from different events may inadequately represent this date. This is a hypothesis suggested by the feature design and saved analysis, not an isolated intervention test. Sunday 3 December is the next-worst day. Neither these diagnoses nor an error recalculated after excluding difficult dates should replace the complete test result.
\begin{table}[tbp]\centering\small
\caption{Five largest daily MAEs for the selected h24 SVR (T08). Bias is prediction minus actual; all numerical columns are in kW.}\label{tab:worst_days}
\setlength{\tabcolsep}{4pt}
\begin{tabular}{lrrrr}
\toprule
Date & MAE & Bias & Actual mean & Predicted mean \\
\midrule
2017-12-01 & \num{4579.2} & \num{-3616.7} & \num{58574.2} & \num{54957.4} \\
2017-12-03 & \num{3028.5} & \num{-2880.5} & \num{59192.4} & \num{56311.9} \\
2017-12-12 & \num{1714.1} & \num{1712.5} & \num{64338.5} & \num{66051.1} \\
2017-12-21 & \num{1336.2} & \num{1118.1} & \num{64373.9} & \num{65492.0} \\
2017-12-02 & \num{1273.4} & \num{-554.7} & \num{60104.3} & \num{59549.7} \\
\bottomrule\end{tabular}
\end{table}


\subsection{Interpretability of the CV-preferred tree model}
The explanatory analysis uses residual RF, the tree candidate with the lowest development mean CV RMSE, rather than the selected SVR. \Cref{fig:importance,fig:shap} therefore describe a different estimator. Built-in importance measures impurity reduction within the residual forest; permutation importance measures the RMSE increase of the complete base-plus-correction predictor after shuffling a test column. Ten repeats with seed 42 are used in the full saved experiment.
\begin{figure}[tbp]\centering
\includegraphics[width=\linewidth]{F19a_feature_importance.png}
\caption{Residual-RF built-in and test permutation importance (F19a). Built-in importance concerns the correction to lag-24; permutation importance concerns the full forecast. Shuffling lag-24 also changes the explicit forecast base, so the two panels measure different quantities.}\label{fig:importance}
\end{figure}
\begin{figure}[tbp]\centering
\includegraphics[width=.88\linewidth]{F19b_shap_summary.png}
\caption{Tree SHAP summary for the residual RF on December inputs (F19b). Contributions are in kW of the learned correction, not the complete load or the selected SVR output.}\label{fig:shap}
\end{figure}

SHAP decomposes model predictions into feature contributions\cite{lundberg2017}. Here the tree explainer sees only the residual ensemble, so it excludes the separate additive base $y_{t-24}$. Weekday encodings contribute prominently to the correction, while the large lag-24 permutation effect partly reflects its direct role in the complete prediction. Correlated lags and redundant calendar encodings can yield unrealistic combinations when independently permuted. These post-test diagnostics explain predictive behavior conditionally; they establish neither causal effects nor selected-SVR feature importance. No numerical SHAP-value table is archived. Additional scatter and dependence plots are retained in \cref{sec:supp_figures}.
\evidence{Notebook 07 cells 4--35; T08 final-test, prediction and worst-day files; T07/T10 selection scores; F17--F19; \code{docs/progress.md} and \code{docs/result.md}, Section 11, for test reruns.}
